\question{Has anonymisation been demonstrated by test, as EDPB Opinion 28/2024 requires?}
\label{sec:anontest}
\statusMeasured

\begin{shortanswer}
Partly, and the honest account has three parts. \textbf{A membership-inference testbed now exists}
and finds no usable signal against the released model at any privacy setting --- including with no
differential privacy at all, where the attacks score \MIANoDPThreshold{} and \MIANoDPShadow{} against
a chance baseline of 0.500. \textbf{The attacks demonstrably work}: the same code reaches
\MIAControlShadow{} against a model deliberately trained to memorise. But \textbf{this is one of the
three attack families the Opinion names}; model inversion and reconstruction are not built, the
result is on synthetic data, and it is therefore evidence about the method rather than a
demonstration of anonymity for the pilot model.
\end{shortanswer}

\subsection{What was built}

EDPB Opinion 28/2024 \cite{edpb2024ai} declines to accept a model as anonymous merely because
differential privacy is configured; it asks for a case-by-case assessment of extraction and
identification risk, naming membership inference, model inversion and reconstruction as the relevant
tests. Two membership-inference attacks are implemented against the released global model:

\begin{itemize}
  \item \textbf{Loss threshold} \cite{yeom2018privacy}. A model fits its training data slightly
        better than unseen data, so per-record loss is itself a membership signal. No attacker
        training required; the standard first-line test.
  \item \textbf{Shadow models} \cite{shokri2017membership}. The attacker trains sixteen models on
        resampled halves of the same distribution, learns how in-set and out-of-set records differ,
        and calibrates each target record's loss against that distribution. Stronger, because it
        corrects for records that are simply easy or hard.
\end{itemize}

Both are scored as attack AUROC. The framing is deliberately attacker-favourable: members are every
practice's training rows and non-members are their held-out rows, so both come from the same
practices and the same distribution and any separation is memorisation rather than population shift.

The pass criterion is one-sided and on effect size --- the 95\% interval's upper bound must stay
within \MIAMargin{} of chance. Significance alone would be the wrong test: with small seed variance
it would flag an attack scoring \emph{below} chance --- performing worse than a coin flip --- as a
failure.

\begin{table}[htbp]
  \centering
  \caption{Membership inference against the released model, \NumSeeds{} seeds, mean $\pm$ s.d.
  Chance is 0.500.}
  \begin{tabular}{llrrl}
    \toprule
    Setting & $\varepsilon$ composed & Threshold attack & Shadow attack & Verdict \\
    \midrule
    \emph{Positive control} (n=40, overfit) & --- & \emph{0.655} & \emph{0.659} & \emph{fires} \\
    \midrule
    No DP at all & --- & 0.501 $\pm$ 0.006 & 0.502 $\pm$ 0.015 & pass \\
    Central DP $\sigma$=0.1 & 1211.8 & 0.502 $\pm$ 0.006 & 0.506 $\pm$ 0.010 & pass \\
    Central DP $\sigma$=0.25 & 244.0 & 0.502 $\pm$ 0.007 & 0.509 $\pm$ 0.010 & pass \\
    Central DP $\sigma$=0.5 & 81.1 & 0.502 $\pm$ 0.006 & 0.502 $\pm$ 0.016 & pass \\
    Central DP $\sigma$=1 & 30.1 & 0.503 $\pm$ 0.008 & 0.506 $\pm$ 0.013 & pass \\
    Central DP $\sigma$=2 & 12.3 & 0.504 $\pm$ 0.011 & 0.505 $\pm$ 0.011 & pass \\
    Local DP $\varepsilon$=50/round & 1283.4 & 0.501 $\pm$ 0.006 & 0.505 $\pm$ 0.012 & pass \\
    Local DP $\varepsilon$=20/round & 257.6 & 0.503 $\pm$ 0.009 & 0.506 $\pm$ 0.009 & pass \\
    Local DP $\varepsilon$=10/round & 85.0 & 0.501 $\pm$ 0.009 & 0.502 $\pm$ 0.011 & pass \\
    Local DP $\varepsilon$=5/round & 31.4 & 0.493 $\pm$ 0.010 & 0.494 $\pm$ 0.008 & pass \\
    Local DP $\varepsilon$=1/round & 4.3 & 0.492 $\pm$ 0.017 & 0.497 $\pm$ 0.022 & pass \\
    \bottomrule
\end{tabular}

  \label{tab:mia}
\end{table}

\begin{figure}[htbp]
  \centering
  \includegraphics[width=\textwidth]{mia}
  \caption{Attack strength across every privacy setting the project can deploy, against the
  positive control. The shaded band is the negligible-advantage margin.}
  \label{fig:mia}
\end{figure}

\subsection{Why the result is what it is --- and what it is not}

The attacks find nothing, and they find nothing \emph{even with no differential privacy applied}.
That deserves a mechanistic explanation rather than a celebration, because the explanation is also
the limitation.

\textbf{The released model has almost no capacity to memorise.} It is logistic regression with
eleven parameters fitted over \NumPatients{} patients. A convex model of that size has nowhere to
store an individual record: its optimum is a property of the dataset as a whole. This is the same
argument made in \S\ref{sec:payload}, now tested rather than asserted.

Two consequences follow, and they point in opposite directions.

\emph{In favour:} it is a genuine and useful finding that the disclosure risk of this model class,
at this scale, is not detectable by the standard attacks. It also means \textbf{differential privacy
is not what is providing the protection here} --- capacity is --- which is directly relevant to
\S\ref{sec:noise}'s finding that DP at a meaningful budget is unaffordable. The project is not
choosing between privacy and utility so much as discovering that the model class was never the
exposure.

\emph{Against, and this must not be lost:}

\begin{enumerate}
  \item \textbf{One of three families.} Model inversion and reconstruction are not implemented. A
        pass on membership inference is not a pass on the Opinion's test.
  \item \textbf{Synthetic data, \NumFeatures{} features, no eGFR.} Real practice data is richer, and
        the strongest CKD predictor is absent here (\S\ref{sec:realdata}). A richer feature space
        gives a model more to memorise. \textbf{The audit must be re-run on real data before any
        anonymity claim is made.}
  \item \textbf{Two attacks, not the strongest available.} Stronger likelihood-ratio attacks
        \cite{carlini2022lira} exist and would be the natural next step.
  \item \textbf{A pass is a necessary condition, never a sufficient one.} It is evidence that can
        support an anonymity argument; it does not constitute one.
\end{enumerate}

\subsection{Why the positive control is part of the result}

An attack that never fires proves nothing, and a report claiming ``no leakage detected'' without
showing the detector works should not be believed. The same two attacks, run against a model
deliberately trained to memorise --- a \MIAControlN{}-patient cohort, many passes, no federated
averaging to wash out the individual fit --- reach \MIAControlThreshold{} and \MIAControlShadow{}.
The testbed detects leakage when leakage is present. That is what licenses reading the near-chance
rows as a finding rather than as a null instrument.
